\documentclass[12pt,a4paper]{article}

% ------------------------------------------------
% Pacotes Fundamentais
% ------------------------------------------------
\usepackage[utf8]{inputenc}
\usepackage[T1]{fontenc}
\usepackage[brazil]{babel}
\usepackage{graphicx}
\usepackage[table]{xcolor}
\usepackage{array}
\usepackage{tabularx}
\usepackage{geometry}
\usepackage{fancyhdr}
\usepackage{mathptmx} % Times New Roman
\usepackage{enumitem} 
\usepackage{amsfonts}
\usepackage{amsmath}
\usepackage{amssymb}
\usepackage{tikz}
\usetikzlibrary{positioning, shapes.geometric, calc}

% Cores Institucionais
\definecolor{maincolor}{RGB}{0,100,50}
\definecolor{forestgreen}{RGB}{0,100,50}

% ------------------------------------------------
% Estilos TikZ para Diagramas ER, Lógicos e Wireframes CRUD
% ------------------------------------------------
\tikzset{
    block/.style={rectangle, draw=forestgreen, thick, minimum width=2.5cm, minimum height=0.9cm, fill=white, align=center, font=\small\bfseries},
    diamond_node/.style={diamond, draw=forestgreen, thick, minimum width=1.8cm, minimum height=0.9cm, fill=white, aspect=2, align=center, font=\small},
    attribute/.style={ellipse, draw=forestgreen, thin, minimum width=1.6cm, minimum height=0.6cm, fill=white, align=center, font=\scriptsize},
    line/.style={draw, thick, -},
    table_block/.style={
        rectangle, 
        draw=forestgreen, 
        thick, 
        rounded corners=4pt, 
        minimum width=4.5cm, 
        fill=green!2, 
        align=left, 
        font=\small\ttfamily, 
        inner sep=8pt
    },
    screen/.style={draw=forestgreen, thick, rectangle, rounded corners=8pt, minimum width=8.5cm, minimum height=7.0cm, fill=white},
    screen_title/.style={fill=forestgreen, text=white, minimum width=8.3cm, minimum height=0.6cm, font=\footnotesize\bfseries\sffamily, align=center},
    btn/.style={draw=forestgreen, fill=green!10, rounded corners=3pt, font=\tiny\sffamily, inner sep=3pt},
    btn_danger/.style={draw=red, fill=red!10, rounded corners=3pt, font=\tiny\sffamily, inner sep=3pt, text=red},
    input_field/.style={draw=gray!50, fill=gray!2, rectangle, minimum width=7.5cm, minimum height=0.55cm, font=\tiny\sffamily, align=left, inner sep=4pt}
}

% ------------------------------------------------ 
% Configuração de Página e Margens
% ------------------------------------------------
\geometry{
    a4paper,
    left=2cm,
    right=2cm,
    top=2cm,
    bottom=2cm,
    headheight=15pt
}

\pagestyle{fancy}
\fancyhf{}
\renewcommand{\headrulewidth}{0pt}
\setlength{\parindent}{0pt}

% ------------------------------------------------
% Macros Customizadas (Caixa de Resposta e Alinhamento)
% ------------------------------------------------
\newcommand{\resposta}[1]{}

\newcolumntype{L}[1]{>{\raggedright\arraybackslash}p{#1}}
\newcolumntype{C}[1]{>{\centering\arraybackslash}p{#1}}

% ------------------------------------------------
% Configurações de Caixas Decorativas (tcolorbox)
% ------------------------------------------------
\usepackage[most]{tcolorbox}
\newtcolorbox{blockbox}[1]{%
  colback=green!5,
  colframe=forestgreen,
  coltitle=white,
  title=#1,
  fonttitle=\bfseries,
  boxrule=1.5pt,
  arc=3mm,
  left=8pt,
  right=8pt,
  top=8pt,
  bottom=8pt
}

% ------------------------------------------------
% Configurações de Código (Listings SQL)
% ------------------------------------------------
\usepackage{listings}
\definecolor{codegray}{rgb}{0.5,0.5,0.5}
\definecolor{codepurple}{rgb}{0.58,0,0.82}
\definecolor{backcolour}{rgb}{0.96,0.96,0.94}

\lstdefinestyle{sqlstyle}{
    backgroundcolor=\color{backcolour},   
    commentstyle=\color{codegray},
    keywordstyle=\color{forestgreen}\bfseries,
    numberstyle=\tiny\color{codegray},
    stringstyle=\color{codepurple},
    basicstyle=\ttfamily\footnotesize,
    breakatwhitespace=false,         
    breaklines=true,                 
    captionpos=b,                    
    keepspaces=true,                 
    numbers=none,                    
    showspaces=false,                
    showstringspaces=false,
    showtabs=false,                  
    tabsize=2,
    language=SQL,
    morekeywords={REFERENCES, ON, CASCADE, SET, NULL, NO, ACTION, RESTRICT, OFFSET, LIMIT, DISTINCT},
    literate={á}{{\'a}}1 {ã}{{\~a}}1 {â}{{\^a}}1 {à}{{\`a}}1 {é}{{\'e}}1 {ê}{{\^e}}1 {í}{{\'i}}1 {ó}{{\'o}}1 {õ}{{\~o}}1 {ô}{{\^o}}1 {ú}{{\'u}}1 {ç}{{\c{c}}}1 {Á}{{\'A}}1 {Ã}{{\~A}}1 {Â}{{\^A}}1 {À}{{\`A}}1 {É}{{\'E}}1 {Ê}{{\^E}}1 {Í}{{\'I}}1 {Ó}{{\'O}}1 {Õ}{{\~O}}1 {Ô}{{\^O}}1 {Ú}{{\'U}}1 {Ç}{{\c{C}}}1
}
\lstset{style=sqlstyle}

% ================================================
\begin{document}
% ================================================

% ------------------------------------------------
% Cabeçalho Institucional Oficial da UniRV
% ------------------------------------------------
\begin{tabularx}{\textwidth}{@{} L{4.5cm} X @{}}
    \raisebox{-0.5cm}{\includegraphics[width=4.5cm]{Universidade_de_Rio_Verde.jpg}} &
    \hfill
    \begin{tabular}{|>{\centering\arraybackslash}m{11.6cm}|}
        \hline
        \cellcolor{forestgreen}\color{white}\rule{0pt}{1.2cm}\textbf{TRABALHO DA N3} \tabularnewline
        \hline
    \end{tabular} 
\end{tabularx}
\vspace{0.3cm}

\begin{tabularx}{\textwidth}{|X|L{2.5cm}|}
\hline
\textbf{Curso} & \textbf{Campus} \\
\textit{Engenharia de Software} & \textit{Rio Verde} \\
\hline
\multicolumn{2}{|l|}{\textbf{Disciplina}} \\  
\multicolumn{2}{|l|}{ \textit{Banco de Dados I}} \\
\hline
\textbf{Professor} & \textbf{Turma} \\
\textit{Sergio Souza Novak} & \textit{A / B} \\
\hline
\multicolumn{2}{|l|}{\textbf{Código da Atividade}} \\
\multicolumn{2}{|l|}{\textit{ESW417}} \\
\hline
\end{tabularx}

\vspace{0.5cm}


\textbf{Contexto:} Uma empresa multinacional de tecnologia (TechCorp) necessita implementar um sistema completo (Backend via API REST, Frontend Web e banco de dados obrigatoriamente em \textbf{PostgreSQL}). A regra de negócio determina que:
\begin{itemize}[leftmargin=*, noitemsep]
    \item Cada funcionário pertence a um único departamento, mas um departamento possui múltiplos funcionários (\textbf{N:1}).
    \item Cada funcionário recebe exatamente um crachá de acesso corporativo exclusivo intransferível e um crachá pertence somente ao funcionário alocado (\textbf{1:1}).
    \item Um funcionário pode atuar em múltiplos projetos simultaneamente, assim como um projeto conta com múltiplos funcionários alocados (\textbf{N:N}).
    \item A hierarquia da empresa exige que cada funcionário possa ter um supervisor direto, que também é um funcionário da mesma empresa (\textbf{Autorelacionamento}).
\end{itemize}

O Frontend web consumirá a API REST para prover as listagens paginadas para o RH corporativo.

\vspace{0.3cm}
\begin{center}
    \textbf{Modelo Conceitual (Diagrama Entidade-Relacionamento)} \\[0.2cm]
    \begin{tikzpicture}[node distance=1cm, >=stealth, thick, scale=0.85, every node/.style={transform shape}]
        % Entidades
        \node[block] (func) {Funcionario};
        \node[block, left=5cm of func] (depto) {Departamento};
        \node[block, right=5cm of func] (projeto) {Projeto};
        \node[block, below=5cm of func] (cracha) {Cracha};
        
        % Relacionamentos
        \node[diamond_node, left=1.5cm of func] (pertence) {Pertence};
        \node[diamond_node, right=1.5cm of func] (trabalha) {Atua};
        \node[diamond_node, below=1.5cm of func] (possui) {Possui};
        
        % Ligações Depto <-> Func
        \path[line] (depto) -- (pertence);
        \path[line] (pertence) -- (func);
        
        % Ligações Func <-> Projeto
        \path[line] (func) -- (trabalha);
        \path[line] (trabalha) -- (projeto);
        
        % Ligações Func <-> Cracha
        \path[line] (func) -- (possui);
        \path[line] (possui) -- (cracha);
        
        % Autorelacionamento Supervisor
        \node[diamond_node, above=1cm of func] (superv) {Supervisiona};
        \draw[line] (func.north west) -- (superv.south west);
        \draw[line] (superv.south east) -- (func.north east);
        
        % Cardinalidades
        \node[above, font=\small\bfseries] at ($(depto.east)!0.3!(pertence.west)$) {1};
        \node[above, font=\small\bfseries] at ($(func.west)!0.3!(pertence.east)$) {N};
        
        \node[above, font=\small\bfseries] at ($(func.east)!0.3!(trabalha.west)$) {N};
        \node[above, font=\small\bfseries] at ($(projeto.west)!0.3!(trabalha.east)$) {N};
        
        \node[left, font=\small\bfseries] at ($(func.south)!0.3!(possui.north)$) {1};
        \node[left, font=\small\bfseries] at ($(cracha.north)!0.3!(possui.south)$) {1};
    \end{tikzpicture}
\end{center}

\begin{center}
    \textbf{Modelo Lógico (Relacional)} \\[0.2cm]
    \begin{tikzpicture}[node distance=2.5cm, >=stealth, thick]
        % Tabelas
        \node[table_block] (depto_tbl) at (0,0) {
            \textbf{tbDepartamento} \\
            \rule{4.1cm}{0.4pt} \\
            \underline{id\_depto} (PK) \\
            nome\_depto: TEXT \\
            orcamento: REAL
        };
        
        \node[table_block] (func_tbl) at (6,0) {
            \textbf{tbFuncionario} \\
            \rule{4.1cm}{0.4pt} \\
            \underline{id\_func} (PK) \\
            nome: TEXT \\
            cargo: TEXT \\
            id\_depto (FK) \\
            id\_supervisor (FK)
        };
        
        \node[table_block] (proj_tbl) at (12,0) {
            \textbf{tbProjeto} \\
            \rule{4.1cm}{0.4pt} \\
            \underline{id\_projeto} (PK) \\
            nome\_projeto: TEXT \\
            data\_inicio: TEXT
        };
        
        \node[table_block] (cracha_tbl) at (3,-4.5) {
            \textbf{tbCracha} \\
            \rule{4.1cm}{0.4pt} \\
            \underline{id\_cracha} (PK) \\
            codigo\_acesso: TEXT \\
            nivel\_permissao: INTEGER \\
            id\_func (FK, UNIQUE)
        };
        
        \node[table_block] (alocacao_tbl) at (9,-4.5) {
            \textbf{tbFuncionarioProjeto} \\
            \rule{4.1cm}{0.4pt} \\
            \underline{id\_func} (PK, FK) \\
            \underline{id\_projeto} (PK, FK) \\
            horas\_semanais: INTEGER
        };
        
        % Linhas FK
        \draw[->, color=forestgreen, ultra thick] (func_tbl.west) -- (depto_tbl.east) 
            node[midway, above, font=\scriptsize\bfseries, text=forestgreen] {FK: id\_depto};
            
        \draw[->, color=forestgreen, ultra thick] (func_tbl.north) .. controls +(1.5,1.5) and +(-1.5,1.5) .. (func_tbl.north)
            node[above, font=\scriptsize\bfseries, text=forestgreen] {FK: id\_supervisor};
            
        \draw[->, color=forestgreen, ultra thick] (cracha_tbl.north) -- (func_tbl.south west)
            node[midway, left, font=\scriptsize\bfseries, text=forestgreen] {FK: id\_func};
            
        \draw[->, color=forestgreen, ultra thick] (alocacao_tbl.north west) -- (func_tbl.south east)
            node[midway, left, font=\scriptsize\bfseries, text=forestgreen] {FK: id\_func};
            
        \draw[->, color=forestgreen, ultra thick] (alocacao_tbl.north east) -- (proj_tbl.south west)
            node[midway, right, font=\scriptsize\bfseries, text=forestgreen] {FK: id\_projeto};
    \end{tikzpicture}
\end{center}

\vspace{0.4cm}

\begin{center}
    \textbf{Protótipos Wireframe do CRUD Completo (Sistema Web)} \\[0.3cm]
    
    % =================================================================
    % 1. DEPARTAMENTO
    % =================================================================
    \textbf{1. Entidade: Departamento} \\[0.2cm]
    
    % --- READ ---
    \resizebox{11cm}{!}{%
    \begin{tikzpicture}[scale=0.95, every node/.style={transform shape}]
      \node[screen] (r_scr) at (0,0) {};
      \node[screen_title, below=0.1cm of r_scr.north] (r_title) {Departamentos - Listagem (R)};
      \node[btn, below=0.2cm of r_title, xshift=-2.2cm] (btn_novo) {+ Novo Departamento};
      
      \node[draw=gray!30, fill=white, minimum width=8.1cm, minimum height=0.8cm, below=0.8cm of r_title] (item1) {};
      \node at (item1.west) [anchor=west, xshift=0.2cm, font=\footnotesize\sffamily] {\textbf{1 - TI} | Orçamento: R\$ 500.000};
      \node[btn, right=0.2cm of item1.east, anchor=east, xshift=-0.6cm] (edit1) {Editar};
      \node[btn_danger, right=0.2cm of item1.east, anchor=east, xshift=-0.2cm] (del1) {x};
      
      \node[draw=gray!30, fill=white, minimum width=8.1cm, minimum height=0.8cm, below=0.15cm of item1] (item2) {};
      \node at (item2.west) [anchor=west, xshift=0.2cm, font=\footnotesize\sffamily] {\textbf{2 - RH} | Orçamento: R\$ 200.000};
      \node[btn, right=0.2cm of item2.east, anchor=east, xshift=-0.6cm] (edit2) {Editar};
      \node[btn_danger, right=0.2cm of item2.east, anchor=east, xshift=-0.2cm] (del2) {x};
      
      \node[draw, fill=forestgreen, text=white, circle, inner sep=2pt, font=\scriptsize, above=0.4cm of r_scr.south, xshift=-1.2cm] (p1) {1};
      \node[draw, circle, inner sep=2pt, font=\scriptsize, right=0.25cm of p1] (p2) {2};
      \node[font=\scriptsize, right=0.25cm of p2] (next) {Próxima >>};
    \end{tikzpicture}%
    }
    \\[0.5cm]
    
    % --- CREATE/UPDATE ---
    \resizebox{11cm}{!}{%
    \begin{tikzpicture}[scale=0.95, every node/.style={transform shape}]
      \node[screen] (cu_scr) at (0,0) {};
      \node[screen_title] at (0, 3.1) {Departamento - Cadastro (C/U)};
      
      \node[font=\footnotesize\bfseries\sffamily, anchor=west] at (-3.75, 2.0) {Nome do Departamento:};
      \node[input_field, anchor=center] at (0, 1.55) {TI};
      
      \node[font=\footnotesize\bfseries\sffamily, anchor=west] at (-3.75, 0.85) {Orçamento Anual (R\$):};
      \node[input_field, anchor=center] at (0, 0.4) {500000.00};
      
      \node[btn, anchor=center, minimum width=2.5cm, minimum height=0.6cm] at (-1.5, -2.8) {Salvar};
      \node[btn_danger, anchor=center, minimum width=2.0cm, minimum height=0.6cm] at (1.3, -2.8) {Cancelar};
    \end{tikzpicture}%
    }
    \\[0.5cm]
    
    % --- DELETE ---
    \resizebox{11cm}{!}{%
    \begin{tikzpicture}[scale=0.95, every node/.style={transform shape}]
      \node[screen] (d_scr) at (0,0) {};
      \node[screen_title, below=0.1cm of d_scr.north] (d_title) {Departamento - Excluir (D)};
      
      \node[draw=red!30, fill=red!5, rounded corners=6pt, minimum width=8.1cm, minimum height=3.5cm, below=0.8cm of d_title] (alert_box) {};
      \node[text=red, font=\footnotesize\bfseries\sffamily, align=center, yshift=1.2cm] at (alert_box) {CONFIRMAÇÃO DE EXCLUSÃO};
      \node[font=\footnotesize\sffamily, align=center, text width=7.5cm, yshift=0.3cm] at (alert_box) {Deseja excluir o Departamento \textbf{TI}?};
      \node[font=\tiny\sffamily, align=center, text=red, text width=7.5cm, yshift=-0.6cm] at (alert_box) {\textbf{Aviso de Integridade:} Não é possível excluir departamentos que ainda possuem funcionários vinculados!};
      
      \node[btn_danger, above=0.35cm of d_scr.south, xshift=-1.35cm, minimum width=2.2cm, minimum height=0.5cm] {Excluir};
      \node[btn, above=0.35cm of d_scr.south, xshift=1.35cm, minimum width=2.2cm, minimum height=0.5cm] {Cancelar};
    \end{tikzpicture}%
    }
    
    \vspace{0.6cm}
    % =================================================================
    % 2. FUNCIONÁRIO
    % =================================================================
    \textbf{2. Entidade: Funcionário} \\[0.2cm]
    
    % --- READ ---
    \resizebox{14cm}{!}{%
    \begin{tikzpicture}[scale=0.95, every node/.style={transform shape}]
      \node[screen] (r_scr) at (0,0) {};
      \node[screen_title, below=0.1cm of r_scr.north] (r_title) {Funcionários - Listagem (R)};
      \node[btn, below=0.2cm of r_title, xshift=-2.8cm] (btn_novo) {+ Novo Funcionário};
      
      \node[draw=gray!30, fill=white, minimum width=8.1cm, minimum height=3.0cm, below=0.8cm of r_title] (tbl_grid) {};
      \node[draw=forestgreen!20, fill=forestgreen!5, minimum width=8.1cm, minimum height=0.6cm, below=0.8cm of r_title] (th) {};
      \node at (th.west) [anchor=west, xshift=0.2cm, font=\tiny\bfseries\sffamily] {Nome \hspace{1.5cm} Cargo \hspace{1.2cm} Depto \hspace{0.5cm} Ações};
      
      \node[draw=gray!10, fill=white, minimum width=8.1cm, minimum height=0.6cm, below=0.0cm of th] (r1) {};
      \node at (r1.west) [anchor=west, xshift=0.2cm, font=\tiny\sffamily] {Ana Souza \hspace{1.1cm} Gerente \hspace{1.0cm} TI \hspace{1.0cm} [Editar] [X]};
      
      \node[draw=gray!10, fill=gray!2, minimum width=8.1cm, minimum height=0.6cm, below=0.0cm of r1] (r2) {};
      \node at (r2.west) [anchor=west, xshift=0.2cm, font=\tiny\sffamily] {Carlos Lima \hspace{0.9cm} Analista \hspace{0.9cm} TI \hspace{1.0cm} [Editar] [X]};
      
      \node[draw, fill=forestgreen, text=white, circle, inner sep=2pt, font=\scriptsize, above=0.4cm of r_scr.south, xshift=-1.2cm] (p1) {1};
      \node[draw, circle, inner sep=2pt, font=\scriptsize, right=0.25cm of p1] (p2) {2};
      \node[font=\scriptsize, right=0.25cm of p2] (next) {Próxima >>};
    \end{tikzpicture}%
    }
    \\[0.5cm]
    
    % --- CREATE/UPDATE ---
    \resizebox{11cm}{!}{%
    \begin{tikzpicture}[scale=0.95, every node/.style={transform shape}]
      \node[screen] (cu_scr) at (0,0) {};
      \node[screen_title] at (0, 3.1) {Funcionário - Cadastro (C/U)};
      
      \node[font=\footnotesize\bfseries\sffamily, anchor=west] at (-3.75, 2.2) {Nome do Funcionário:};
      \node[input_field, anchor=center] at (0, 1.78) {Carlos Lima};
      
      \node[font=\footnotesize\bfseries\sffamily, anchor=west] at (-3.75, 1.1) {Cargo:};
      \node[input_field, anchor=center] at (0, 0.68) {Analista Pleno};
      
      \node[font=\footnotesize\bfseries\sffamily, anchor=west] at (-3.75, 0.0) {Departamento (FK):};
      \node[input_field, anchor=center] at (0, -0.42) {1 - TI};
      
      \node[font=\footnotesize\bfseries\sffamily, anchor=west] at (-3.75, -1.1) {Supervisor (Autorelacionamento):};
      \node[input_field, anchor=center] at (0, -1.52) {1 - Ana Souza (Gerente)};
      
      \node[btn, anchor=center, minimum width=2.5cm, minimum height=0.6cm] at (-1.5, -2.8) {Salvar};
      \node[btn_danger, anchor=center, minimum width=2.0cm, minimum height=0.6cm] at (1.3, -2.8) {Cancelar};
    \end{tikzpicture}%
    }
    \\[0.5cm]
    
    % --- DELETE ---
    \resizebox{11cm}{!}{%
    \begin{tikzpicture}[scale=0.95, every node/.style={transform shape}]
      \node[screen] (d_scr) at (0,0) {};
      \node[screen_title, below=0.1cm of d_scr.north] (d_title) {Funcionário - Excluir (D)};
      
      \node[draw=red!30, fill=red!5, rounded corners=6pt, minimum width=8.1cm, minimum height=3.5cm, below=0.8cm of d_title] (alert_box) {};
      \node[text=red, font=\footnotesize\bfseries\sffamily, align=center, yshift=1.2cm] at (alert_box) {CONFIRMAÇÃO DE EXCLUSÃO};
      \node[font=\footnotesize\sffamily, align=center, text width=7.5cm, yshift=0.3cm] at (alert_box) {Deseja excluir \textbf{Carlos Lima}?};
      \node[font=\tiny\sffamily, align=center, text=red, text width=7.5cm, yshift=-0.6cm] at (alert_box) {\textbf{Aviso de Integridade:} Ao excluir, suas alocações em projetos e seu crachá (CASCADE) serão removidos. Se ele for supervisor, os supervisionados ficarão órfãos!};
      
      \node[btn_danger, above=0.35cm of d_scr.south, xshift=-1.35cm, minimum width=2.2cm, minimum height=0.5cm] {Excluir};
      \node[btn, above=0.35cm of d_scr.south, xshift=1.35cm, minimum width=2.2cm, minimum height=0.5cm] {Cancelar};
    \end{tikzpicture}%
    }
    
    \vspace{0.6cm}
    % =================================================================
    % 3. PROJETO
    % =================================================================
    \textbf{3. Entidade: Projeto} \\[0.2cm]
    
    % --- READ ---
    \resizebox{11cm}{!}{%
    \begin{tikzpicture}[scale=0.95, every node/.style={transform shape}]
      \node[screen] (r_scr) at (0,0) {};
      \node[screen_title, below=0.1cm of r_scr.north] (r_title) {Projetos - Listagem (R)};
      \node[btn, below=0.2cm of r_title, xshift=-2.6cm] (btn_novo) {+ Novo Projeto};
      
      \node[draw=gray!30, fill=white, minimum width=8.1cm, minimum height=0.8cm, below=0.8cm of r_title] (item1) {};
      \node at (item1.west) [anchor=west, xshift=0.2cm, font=\footnotesize\sffamily] {\textbf{Alpha} | Início: 01/01/26};
      \node[btn, right=0.2cm of item1.east, anchor=east, xshift=-0.6cm] (edit1) {Editar};
      \node[btn_danger, right=0.2cm of item1.east, anchor=east, xshift=-0.2cm] (del1) {x};
      
      \node[draw=gray!30, fill=white, minimum width=8.1cm, minimum height=0.8cm, below=0.15cm of item1] (item2) {};
      \node at (item2.west) [anchor=west, xshift=0.2cm, font=\footnotesize\sffamily] {\textbf{Beta} | Início: 15/02/26};
      \node[btn, right=0.2cm of item2.east, anchor=east, xshift=-0.6cm] (edit2) {Editar};
      \node[btn_danger, right=0.2cm of item2.east, anchor=east, xshift=-0.2cm] (del2) {x};
      
      \node[draw, fill=forestgreen, text=white, circle, inner sep=2pt, font=\scriptsize, above=0.4cm of r_scr.south, xshift=-1.2cm] (p1) {1};
      \node[draw, circle, inner sep=2pt, font=\scriptsize, right=0.25cm of p1] (p2) {2};
      \node[font=\scriptsize, right=0.25cm of p2] (next) {Próxima >>};
    \end{tikzpicture}%
    }
    \\[0.5cm]
    
    % --- CREATE/UPDATE ---
    \resizebox{11cm}{!}{%
    \begin{tikzpicture}[scale=0.95, every node/.style={transform shape}]
      \node[screen] (cu_scr) at (0,0) {};
      \node[screen_title] at (0, 3.1) {Projeto - Cadastro (C/U)};
      
      \node[font=\footnotesize\bfseries\sffamily, anchor=west] at (-3.75, 2.0) {Nome do Projeto:};
      \node[input_field, anchor=center] at (0, 1.55) {Alpha};
      
      \node[font=\footnotesize\bfseries\sffamily, anchor=west] at (-3.75, 0.85) {Data de Início:};
      \node[input_field, anchor=center] at (0, 0.4) {01/01/2026};
      
      \node[btn, anchor=center, minimum width=2.5cm, minimum height=0.6cm] at (-1.5, -2.8) {Salvar};
      \node[btn_danger, anchor=center, minimum width=2.0cm, minimum height=0.6cm] at (1.3, -2.8) {Cancelar};
    \end{tikzpicture}%
    }
    \\[0.5cm]
    
    % --- DELETE ---
    \resizebox{11cm}{!}{%
    \begin{tikzpicture}[scale=0.95, every node/.style={transform shape}]
      \node[screen] (d_scr) at (0,0) {};
      \node[screen_title, below=0.1cm of d_scr.north] (d_title) {Projeto - Excluir (D)};
      
      \node[draw=red!30, fill=red!5, rounded corners=6pt, minimum width=8.1cm, minimum height=3.5cm, below=0.8cm of d_title] (alert_box) {};
      \node[text=red, font=\footnotesize\bfseries\sffamily, align=center, yshift=1.2cm] at (alert_box) {CONFIRMAÇÃO DE EXCLUSÃO};
      \node[font=\footnotesize\sffamily, align=center, text width=7.5cm, yshift=0.3cm] at (alert_box) {Deseja excluir o Projeto \textbf{Alpha}?};
      \node[font=\tiny\sffamily, align=center, text=red, text width=7.5cm, yshift=-0.6cm] at (alert_box) {\textbf{Aviso de Integridade:} Isso removerá os vínculos na tabela N:N (FuncionarioProjeto).};
      
      \node[btn_danger, above=0.35cm of d_scr.south, xshift=-1.35cm, minimum width=2.2cm, minimum height=0.5cm] {Excluir};
      \node[btn, above=0.35cm of d_scr.south, xshift=1.35cm, minimum width=2.2cm, minimum height=0.5cm] {Cancelar};
    \end{tikzpicture}%
    }
    
    \vspace{0.6cm}
    % =================================================================
    % 4. CRACHÁ
    % =================================================================
    \textbf{4. Entidade: Crachá (1:1)} \\[0.2cm]
    
    % --- READ ---
    \resizebox{11cm}{!}{%
    \begin{tikzpicture}[scale=0.95, every node/.style={transform shape}]
      \node[screen] (r_scr) at (0,0) {};
      \node[screen_title, below=0.1cm of r_scr.north] (r_title) {Crachás - Listagem (R)};
      \node[btn, below=0.2cm of r_title, xshift=-2.6cm] (btn_novo) {+ Novo Crachá};
      
      \node[draw=gray!30, fill=white, minimum width=8.1cm, minimum height=0.8cm, below=0.8cm of r_title] (item1) {};
      \node at (item1.west) [anchor=west, xshift=0.2cm, font=\footnotesize\sffamily] {\textbf{X99} | Func: Ana | Nível: 5};
      \node[btn, right=0.2cm of item1.east, anchor=east, xshift=-0.6cm] (edit1) {Editar};
      \node[btn_danger, right=0.2cm of item1.east, anchor=east, xshift=-0.2cm] (del1) {x};
      
      \node[draw, fill=forestgreen, text=white, circle, inner sep=2pt, font=\scriptsize, above=0.4cm of r_scr.south, xshift=-1.2cm] (p1) {1};
      \node[draw, circle, inner sep=2pt, font=\scriptsize, right=0.25cm of p1] (p2) {2};
      \node[font=\scriptsize, right=0.25cm of p2] (next) {Próxima >>};
    \end{tikzpicture}%
    }
    \\[0.5cm]
    
    % --- CREATE/UPDATE ---
    \resizebox{11cm}{!}{%
    \begin{tikzpicture}[scale=0.95, every node/.style={transform shape}]
      \node[screen] (cu_scr) at (0,0) {};
      \node[screen_title] at (0, 3.1) {Crachá - Cadastro (C/U)};
      
      \node[font=\footnotesize\bfseries\sffamily, anchor=west] at (-3.75, 2.0) {Código de Acesso:};
      \node[input_field, anchor=center] at (0, 1.55) {X99};
      
      \node[font=\footnotesize\bfseries\sffamily, anchor=west] at (-3.75, 0.85) {Nível de Permissão:};
      \node[input_field, anchor=center] at (0, 0.4) {5};
      
      \node[font=\footnotesize\bfseries\sffamily, anchor=west] at (-3.75, -0.3) {Funcionário Vinculado (FK):};
      \node[input_field, anchor=center] at (0, -0.75) {1 - Ana Souza};
      
      \node[btn, anchor=center, minimum width=2.5cm, minimum height=0.6cm] at (-1.5, -2.8) {Salvar};
      \node[btn_danger, anchor=center, minimum width=2.0cm, minimum height=0.6cm] at (1.3, -2.8) {Cancelar};
    \end{tikzpicture}%
    }
    \\[0.5cm]
    
    % --- DELETE ---
    \resizebox{11cm}{!}{%
    \begin{tikzpicture}[scale=0.95, every node/.style={transform shape}]
      \node[screen] (d_scr) at (0,0) {};
      \node[screen_title, below=0.1cm of d_scr.north] (d_title) {Crachá - Inativar (D)};
      
      \node[draw=red!30, fill=red!5, rounded corners=6pt, minimum width=8.1cm, minimum height=3.5cm, below=0.8cm of d_title] (alert_box) {};
      \node[text=red, font=\footnotesize\bfseries\sffamily, align=center, yshift=1.2cm] at (alert_box) {CONFIRMAÇÃO DE INATIVAÇÃO};
      \node[font=\footnotesize\sffamily, align=center, text width=7.5cm, yshift=0.3cm] at (alert_box) {Deseja inutilizar o crachá \textbf{X99}?};
      \node[font=\tiny\sffamily, align=center, text=red, text width=7.5cm, yshift=-0.6cm] at (alert_box) {\textbf{Aviso de Integridade:} O funcionário Ana Souza perderá o acesso físico temporariamente até novo vínculo.};
      
      \node[btn_danger, above=0.35cm of d_scr.south, xshift=-1.35cm, minimum width=2.2cm, minimum height=0.5cm] {Excluir};
      \node[btn, above=0.35cm of d_scr.south, xshift=1.35cm, minimum width=2.2cm, minimum height=0.5cm] {Cancelar};
    \end{tikzpicture}%
    }
    
    \vspace{0.6cm}
    % =================================================================
    % 5. ALOCAÇÃO (FuncionarioProjeto)
    % =================================================================
    \textbf{5. Entidade Associativa: Alocação (N:N)} \\[0.2cm]
    
    % --- READ ---
    \resizebox{14cm}{!}{%
    \begin{tikzpicture}[scale=0.95, every node/.style={transform shape}]
      \node[screen] (r_scr) at (0,0) {};
      \node[screen_title, below=0.1cm of r_scr.north] (r_title) {Alocações de Projetos - Listagem (R)};
      \node[btn, below=0.2cm of r_title, xshift=-2.8cm] (btn_novo) {+ Nova Alocação};
      
      \node[draw=gray!30, fill=white, minimum width=8.1cm, minimum height=3.0cm, below=0.8cm of r_title] (tbl_grid) {};
      \node[draw=forestgreen!20, fill=forestgreen!5, minimum width=8.1cm, minimum height=0.6cm, below=0.8cm of r_title] (th) {};
      \node at (th.west) [anchor=west, xshift=0.2cm, font=\tiny\bfseries\sffamily] {Funcionário \hspace{1.0cm} Projeto \hspace{1.0cm} Horas \hspace{0.5cm} Ações};
      
      \node[draw=gray!10, fill=white, minimum width=8.1cm, minimum height=0.6cm, below=0.0cm of th] (r1) {};
      \node at (r1.west) [anchor=west, xshift=0.2cm, font=\tiny\sffamily] {Carlos Lima \hspace{1.1cm} Alpha \hspace{1.3cm} 20h \hspace{1.0cm} [Editar] [X]};
      
      \node[draw=gray!10, fill=gray!2, minimum width=8.1cm, minimum height=0.6cm, below=0.0cm of r1] (r2) {};
      \node at (r2.west) [anchor=west, xshift=0.2cm, font=\tiny\sffamily] {Carlos Lima \hspace{1.1cm} Beta \hspace{1.4cm} 20h \hspace{1.0cm} [Editar] [X]};
      
      \node[draw, fill=forestgreen, text=white, circle, inner sep=2pt, font=\scriptsize, above=0.4cm of r_scr.south, xshift=-1.2cm] (p1) {1};
      \node[draw, circle, inner sep=2pt, font=\scriptsize, right=0.25cm of p1] (p2) {2};
      \node[font=\scriptsize, right=0.25cm of p2] (next) {Próxima >>};
    \end{tikzpicture}%
    }
    \\[0.5cm]
    
    % --- CREATE/UPDATE ---
    \resizebox{11cm}{!}{%
    \begin{tikzpicture}[scale=0.95, every node/.style={transform shape}]
      \node[screen] (cu_scr) at (0,0) {};
      \node[screen_title] at (0, 3.1) {Alocação - Cadastro (C/U)};
      
      \node[font=\footnotesize\bfseries\sffamily, anchor=west] at (-3.75, 2.0) {Funcionário (FK id\_func):};
      \node[input_field, anchor=center] at (0, 1.55) {2 - Carlos Lima};
      
      \node[font=\footnotesize\bfseries\sffamily, anchor=west] at (-3.75, 0.85) {Projeto (FK id\_projeto):};
      \node[input_field, anchor=center] at (0, 0.4) {1 - Alpha};
      
      \node[font=\footnotesize\bfseries\sffamily, anchor=west] at (-3.75, -0.3) {Horas Semanais:};
      \node[input_field, anchor=center] at (0, -0.75) {20};
      
      \node[btn, anchor=center, minimum width=2.5cm, minimum height=0.6cm] at (-1.5, -2.8) {Salvar};
      \node[btn_danger, anchor=center, minimum width=2.0cm, minimum height=0.6cm] at (1.3, -2.8) {Cancelar};
    \end{tikzpicture}%
    }
    \\[0.5cm]
    
    % --- DELETE ---
    \resizebox{11cm}{!}{%
    \begin{tikzpicture}[scale=0.95, every node/.style={transform shape}]
      \node[screen] (d_scr) at (0,0) {};
      \node[screen_title, below=0.1cm of d_scr.north] (d_title) {Alocação - Desvincular (D)};
      
      \node[draw=red!30, fill=red!5, rounded corners=6pt, minimum width=8.1cm, minimum height=3.5cm, below=0.8cm of d_title] (alert_box) {};
      \node[text=red, font=\footnotesize\bfseries\sffamily, align=center, yshift=1.2cm] at (alert_box) {CONFIRMAÇÃO DE DESVÍNCULO};
      \node[font=\footnotesize\sffamily, align=center, text width=7.5cm, yshift=0.3cm] at (alert_box) {Deseja desalocar \textbf{Carlos Lima} do Projeto \textbf{Alpha}?};
      \node[font=\tiny\sffamily, align=center, text=red, text width=7.5cm, yshift=-0.6cm] at (alert_box) {\textbf{Aviso de Integridade:} Essa ação exclui o relacionamento. Funcionário e Projeto permanecem intactos no banco.};
      
      \node[btn_danger, above=0.35cm of d_scr.south, xshift=-1.35cm, minimum width=2.2cm, minimum height=0.5cm] {Desvincular};
      \node[btn, above=0.35cm of d_scr.south, xshift=1.35cm, minimum width=2.2cm, minimum height=0.5cm] {Cancelar};
    \end{tikzpicture}%
    }
\end{center}


\vspace{0.4cm}
\section*{Enunciado do Projeto Final}

O objetivo deste trabalho prático é a \textbf{implementação completa (Full-Stack)} do sistema de gestão corporativa da TechCorp, baseando-se nos modelos conceitual, lógico e nos protótipos de interface (wireframes) fornecidos nas páginas anteriores. 

Os requisitos obrigatórios para a aprovação são:
\begin{itemize}[leftmargin=*, noitemsep]
    \item \textbf{Banco de Dados (Obrigatório PostgreSQL):} A estrutura física do banco deve ser implementada de forma relacional utilizando obrigatoriamente o PostgreSQL. O aluno deverá escrever e entregar o script DDL contendo a criação das tabelas, definição das chaves primárias e as restrições de chaves estrangeiras apropriadas (incluindo tratamento de exclusões em \texttt{CASCADE} ou \texttt{RESTRICT} conforme os alertas nos wireframes).
    \item \textbf{Backend (API REST):} Desenvolvimento de um servidor backend (em linguagem de escolha livre: Node.js, Python FastAPI, Java Spring, C\#, PHP, etc.) que conecte ao PostgreSQL e proveja as rotas (endpoints) para as operações de CRUD. As consultas de listagem na API devem, obrigatoriamente, suportar \textbf{Paginação} no banco via parâmetros na rota (utilizando as cláusulas \texttt{LIMIT} e \texttt{OFFSET}).
    \item \textbf{Frontend Web:} Implementação das telas e navegação do sistema web. A interface deve ser fiel à arquitetura de informações e fluxos propostos nos protótipos wireframe deste documento. O Frontend deve consumir a API REST assincronamente e tratar os retornos e erros do sistema de forma adequada para o usuário.
\end{itemize}

\vspace{0.4cm}
\begin{blockbox}{Critérios e Etapas de Entrega}
O trabalho prático pode ser realizado individualmente ou em \textbf{duplas}. A avaliação será dividida obrigatoriamente em duas etapas:

\textbf{Etapa 1: Submissão Digital (Google Classroom)}
\begin{itemize}[leftmargin=*, noitemsep]
    \item \textbf{Ambos} os integrantes da dupla devem realizar a submissão da atividade no Classroom.
    \item A entrega deve conter todo o código-fonte integrado dos componentes do sistema (Backend API, Frontend Web e script \texttt{.sql} do PostgreSQL).
    \item Um \textbf{vídeo gravado} (com duração de 10 a 15 minutos) demonstrando o sistema em pleno funcionamento. O vídeo deve provar, na prática, o funcionamento das 4 operações de CRUD (Create, Read, Update, Delete) para todas as entidades e exibir a implementação bem-sucedida da paginação de listagens.
\end{itemize}

\textbf{Etapa 2: Apresentação Presencial (UniRV)}
\begin{itemize}[leftmargin=*, noitemsep]
    \item A nota final será validada através da demonstração presencial do projeto executando na carteira do aluno na Universidade de Rio Verde. A avaliação e arguição direta serão realizadas por uma \textbf{banca avaliadora} composta pelo \textbf{Professor} da disciplina e um \textbf{Discente} (aluno da turma) escolhido previamente pelo professor.
\end{itemize}
\end{blockbox}

\newpage
\begin{center}
    {\huge \textbf{\textcolor{forestgreen}{Ficha Trabalho Prático}}} \\[0.3cm]
    {\large \textbf{Avaliação Geral de Implementação e Apresentação}} \\[0.8cm]
\end{center}

Esta ficha de avaliação oficial balizará o cálculo da nota final do acadêmico, cuja composição considera tanto a qualidade do código entregue (itens 1 a 7) quanto o desempenho individual durante a arguição (item 8). A lógica avaliativa estabelece que a nota base do projeto (o somatório obtido nos itens de 1 a 7) terá 60\% de seu valor total garantido mediante a correta implementação e submissão do sistema no formato estipulado.

Os 40\% restantes dessa nota base ficam expressamente condicionados ao desempenho prático do acadêmico diante da banca. Durante a arguição, será atribuída uma nota individual de 0 a 100 que atuará como rendimento ponderado sobre essa fatia de 40\% do projeto. Em suma, o sistema aplica a fórmula matemática: \textbf{Nota Final = Nota Base $\times$ 0,6 + $\Big($Nota Base $\times$ 0,4 $\times \frac{\text{Nota da Banca}}{100}\Big)$}. Esse formato rigoroso assegura que a pontuação do trabalho só seja integralizada caso o aluno comprove pleno domínio arquitetural e lógico sobre o código-fonte desenvolvido.
\vspace{0.5cm}

\begin{center}
    
    \renewcommand{\arraystretch}{1.5}
    \begin{tabularx}{\textwidth}{|c|X|c|c|}
        \hline
        \rowcolor{forestgreen}\textcolor{white}{\textbf{Item}} & \textcolor{white}{\textbf{Critério de Avaliação}} & \textcolor{white}{\textbf{Peso}} & \textcolor{white}{\textbf{Nota}} \\
        \hline
        1 & \textbf{Banco de Dados - Estrutura Física:} Scripts DDL consistentes, tabelas e chaves primárias/estrangeiras (PK/FK) em PostgreSQL. & 2.5 & \\
        \hline
        2 & \textbf{Banco de Dados - Integridade:} Definição correta das regras de exclusão (CASCADE/RESTRICT) e validação do autorelacionamento. & 2.5 & \\
        \hline
        3 & \textbf{Backend - Arquitetura e Rotas:} Implementação da API REST com rotas organizadas e métodos HTTP corretos (GET, POST, PUT, DELETE). & 1.5 & \\
        \hline
        4 & \textbf{Backend - Paginação:} Lógica SQL via API garantindo a aplicação de \texttt{LIMIT/OFFSET} para listagens paginadas com performance. & 1.5 & \\
        \hline
        5 & \textbf{Frontend - UI e UX:} Interface amigável e fidedigna aos protótipos de wireframes contidos no escopo do trabalho. & 0.5 & \\
        \hline
        6 & \textbf{Frontend - Integração REST:} Consumo assíncrono da API com sucesso, preenchimento dinâmico de tabelas e exibição de alertas de erro. & 0.5 & \\
        \hline
        7 & \textbf{Entrega (Classroom e Vídeo):} Envio do código-fonte integrado no prazo mais a demonstração gravada com qualidade (vídeo de 10 a 15 min). & 1.0 & \\
        \hline
        \rowcolor{gray!10}
        8 & \textbf{Arguição Presencial (Nota da Banca):} Média da avaliação sobre o domínio profundo do código-fonte e respostas às perguntas (Professor e Discente convidado). & \textbf{0 a 100} & \\
        \hline
        \multicolumn{2}{|r|}{\textbf{NOTA FINAL: Soma [1 a 7] $\times$ 0,6 + $\Big($Soma [1 a 7] $\times$ 0,4 $\times \frac{\text{Nota Banca [8]}}{100}\Big)$}} & \textbf{10.0} & \\
        \hline
    \end{tabularx}
\end{center}

\vspace{0.4cm}
\begin{center}
    \textcolor{forestgreen}{\textbf{\Large $\star$ Bônus de Desempenho $\star$}} \\[0.2cm]
    \textbf{Atenção:} As \textbf{3 melhores duplas} avaliadas na entrega e arguição deste trabalho receberão uma bonificação de \textbf{15 pontos extras } do total de 100 pontos na média final da avaliação N3!
\end{center}
% =========================================================
% FOLHA 1 - AVALIADOR 1
% =========================================================
\newpage
\thispagestyle{empty}

\begin{center}
    {\huge \textbf{\textcolor{forestgreen}{Ficha de Arguição Presencial}}} \\[0.3cm]
    {\large \textbf{Avaliação Individual de Implementação e Apresentação}} \\[0.8cm]
\end{center}

\begin{tcolorbox}[colback=white, colframe=forestgreen, title=\textbf{Identificação da Avaliação}, fonttitle=\bfseries]
    \renewcommand{\arraystretch}{1.5}
    \begin{tabular}{p{0.95\textwidth}}
         \textbf{Membro da Banca:} [~~] Professor \hspace{1cm} [~~] Discente Convidado \\
         \textbf{Nome do Avaliador:} \underline{\hspace{10cm}} \\
         \textbf{Dupla Avaliada (Tema/Projeto):} \underline{\hspace{8.5cm}} \\
         \textbf{Acadêmico 1 (A1):} \underline{\hspace{12cm}} \\
         \textbf{Acadêmico 2 (A2):} \underline{\hspace{12cm}} \\
    \end{tabular}
\end{tcolorbox}

\vspace{0.8cm}

\begin{center}
    \renewcommand{\arraystretch}{2.2}
    \begin{tabularx}{\textwidth}{|X|c|c|}
        \hline
        \rowcolor{forestgreen}\textcolor{white}{\textbf{Critérios Avaliados na Arguição (Atribua uma nota de 0 a 100)}} & \textcolor{white}{\textbf{Acadêmico 1}} & \textcolor{white}{\textbf{Acadêmico 2}} \\
        \hline
        \multicolumn{3}{|c|}{\cellcolor{gray!20}\textbf{Aspectos de Implementação (Domínio do Código)}} \\
        \hline
        1. \textbf{Banco de Dados:} Demonstrou conhecimento sobre o esquema físico, DDL, PKs, FKs e restrições de integridade desenvolvidos? & & \\
        \hline
        2. \textbf{Backend:} Explicou com clareza o funcionamento da lógica da API REST e como a paginação foi construída no SQL? & & \\
        \hline
        3. \textbf{Frontend e Integração:} Sabe explicar o consumo dos endpoints, a construção da interface e validações tratadas? & & \\
        \hline
        \multicolumn{3}{|c|}{\cellcolor{gray!20}\textbf{Aspectos de Apresentação e Postura}} \\
        \hline
        4. \textbf{Desenvoltura Técnica:} Respondeu aos questionamentos de forma assertiva e com vocabulário técnico adequado? & & \\
        \hline
        5. \textbf{Divisão de Tarefas:} Participação equilibrada na arguição e conhecimento da estrutura geral (full-stack) da aplicação? & & \\
        \hline
        \rowcolor{gray!10}
        \multicolumn{1}{|r|}{\textbf{NOTA FINAL DO ALUNO (0 a 100):}} & & \\
        \hline
    \end{tabularx}
\end{center}

\vspace{2cm}
\begin{center}
    \underline{\hspace{8cm}} \\
    \textbf{Assinatura do Avaliador}
\end{center}

% =========================================================
% FOLHA 2 - AVALIADOR 2
% =========================================================
\newpage
\thispagestyle{empty}

\begin{center}
    {\huge \textbf{\textcolor{forestgreen}{Ficha de Arguição Presencial}}} \\[0.3cm]
    {\large \textbf{Avaliação Individual de Implementação e Apresentação}} \\[0.8cm]
\end{center}

\begin{tcolorbox}[colback=white, colframe=forestgreen, title=\textbf{Identificação da Avaliação}, fonttitle=\bfseries]
    \renewcommand{\arraystretch}{1.5}
    \begin{tabular}{p{0.95\textwidth}}
         \textbf{Membro da Banca:} [~~] Professor \hspace{1cm} [~~] Discente Convidado \\
         \textbf{Nome do Avaliador:} \underline{\hspace{10cm}} \\
         \textbf{Dupla Avaliada (Tema/Projeto):} \underline{\hspace{8.5cm}} \\
         \textbf{Acadêmico 1 (A1):} \underline{\hspace{12cm}} \\
         \textbf{Acadêmico 2 (A2):} \underline{\hspace{12cm}} \\
    \end{tabular}
\end{tcolorbox}

\vspace{0.8cm}

\begin{center}
    \renewcommand{\arraystretch}{2.2}
    \begin{tabularx}{\textwidth}{|X|c|c|}
        \hline
        \rowcolor{forestgreen}\textcolor{white}{\textbf{Critérios Avaliados na Arguição (Atribua o multiplicador de 0.0 a 1.0)}} & \textcolor{white}{\textbf{Acadêmico 1}} & \textcolor{white}{\textbf{Acadêmico 2}} \\
        \hline
        \multicolumn{3}{|c|}{\cellcolor{gray!20}\textbf{Aspectos de Implementação (Domínio do Código)}} \\
        \hline
        1. \textbf{Banco de Dados:} Demonstrou conhecimento sobre o esquema físico, DDL, PKs, FKs e restrições de integridade desenvolvidos? & & \\
        \hline
        2. \textbf{Backend:} Explicou com clareza o funcionamento da lógica da API REST e como a paginação foi construída no SQL? & & \\
        \hline
        3. \textbf{Frontend e Integração:} Sabe explicar o consumo dos endpoints, a construção da interface e validações tratadas? & & \\
        \hline
        \multicolumn{3}{|c|}{\cellcolor{gray!20}\textbf{Aspectos de Apresentação e Postura}} \\
        \hline
        4. \textbf{Desenvoltura Técnica:} Respondeu aos questionamentos de forma assertiva e com vocabulário técnico adequado? & & \\
        \hline
        5. \textbf{Divisão de Tarefas:} Participação equilibrada na arguição e conhecimento da estrutura geral (full-stack) da aplicação? & & \\
        \hline
        \rowcolor{gray!10}
        \multicolumn{1}{|r|}{\textbf{MÉDIA MULTIPLICADORA DO ALUNO (0.0 a 1.0):}} & & \\
        \hline
    \end{tabularx}
\end{center}

\vspace{2cm}
\begin{center}
    \underline{\hspace{8cm}} \\
    \textbf{Assinatura do Avaliador}
\end{center}

\end{document}
